\documentclass[11pt,a4paper]{article}
\usepackage[top=2.3cm,bottom=2.3cm,left=2.2cm,right=2.2cm]{geometry}

\usepackage{amsmath,amssymb}
\usepackage{fontspec}
\setmainfont{Liberation Serif}
\setsansfont{Liberation Sans}
\setmonofont{DejaVu Sans Mono}[Scale=0.84]

\usepackage{xcolor}
\definecolor{brandblue}{RGB}{20, 55, 105}
\definecolor{accentblue}{RGB}{35, 95, 165}
\definecolor{codebg}{RGB}{248, 249, 251}
\definecolor{codeframe}{RGB}{215, 222, 230}
\definecolor{javakeyword}{RGB}{0, 45, 170}
\definecolor{javastring}{RGB}{5, 120, 20}
\definecolor{javacomment}{RGB}{120, 120, 120}
\definecolor{javanumber}{RGB}{20, 75, 220}

\definecolor{noteborder}{RGB}{30, 110, 65}
\definecolor{notebg}{RGB}{242, 249, 245}
\definecolor{warnborder}{RGB}{185, 75, 10}
\definecolor{warnbg}{RGB}{254, 248, 240}
\definecolor{tipborder}{RGB}{30, 95, 160}
\definecolor{tipbg}{RGB}{242, 247, 254}

\usepackage{fancyhdr}
\usepackage{lastpage}
\pagestyle{fancy}
\fancyhf{}
\lhead{\parbox[t]{0.58\textwidth}{\raggedright\small\sffamily\color{brandblue}\textbf{T12 --- Ereditarietà, Dynamic Binding e Polimorfismo}}}
\rhead{\small\sffamily\color{gray}Programmazione II -- UniVR (2025/2026)}
\cfoot{\small\sffamily Pagina \thepage\ di \pageref{LastPage}}
\renewcommand{\headrulewidth}{0.5pt}
\renewcommand{\headrule}{\hbox to\headwidth{\color{brandblue!70!black}\leaders\hrule height \headrulewidth\hfill}}

\usepackage{titlesec}
\titleformat{\section}{\Large\bfseries\sffamily\color{brandblue}}{\thesection}{0.8em}{}[\titlerule]
\titleformat{\subsection}{\large\bfseries\sffamily\color{accentblue}}{\thesubsection}{0.8em}{}
\titleformat{\subsubsection}{\normalsize\bfseries\sffamily\color{brandblue!85!black}}{\thesubsubsection}{0.8em}{}

\usepackage{needspace}
\usepackage{fancyvrb}
\DefineVerbatimEnvironment{asciibox}{Verbatim}{
  fontsize=\fontsize{7.5pt}{9.2pt}\selectfont,
  frame=single,
  rulecolor=\color{codeframe},
  fillcolor=\color{codebg},
  framesep=2.5mm
}
\usepackage{listings}
\lstset{
  backgroundcolor=\color{codebg},
  frame=single,
  rulecolor=\color{codeframe},
  basicstyle=\ttfamily\footnotesize,
  keywordstyle=\color{javakeyword}\bfseries,
  stringstyle=\color{javastring},
  commentstyle=\color{javacomment}\itshape,
  numberstyle=\tiny\color{gray},
  numbers=left,
  numbersep=7pt,
  stepnumber=1,
  breaklines=true,
  breakatwhitespace=true,
  showstringspaces=false,
  tabsize=4,
  xleftmargin=12pt,
  framexleftmargin=12pt
}

\newcommand{\passthrough}[1]{#1}
\providecommand{\tightlist}{\setlength{\itemsep}{0pt}\setlength{\parskip}{0pt}}

\usepackage{tcolorbox}
\tcbuselibrary{skins,breakable}
\newtcolorbox{studynote}[1][Nota]{
  enhanced,
  breakable,
  colback=notebg,
  colframe=noteborder,
  fonttitle=\bfseries\sffamily\small,
  coltitle=white,
  title={#1},
  arc=2mm,
  boxrule=0.8pt,
  left=9pt, right=9pt, top=7pt, bottom=7pt
}

\newtcolorbox{studywarning}[1][Attenzione]{
  enhanced,
  breakable,
  colback=warnbg,
  colframe=warnborder,
  fonttitle=\bfseries\sffamily\small,
  coltitle=white,
  title={#1},
  arc=2mm,
  boxrule=0.8pt,
  left=9pt, right=9pt, top=7pt, bottom=7pt
}

\newtcolorbox{studytip}[1][Suggerimento]{
  enhanced,
  breakable,
  colback=tipbg,
  colframe=tipborder,
  fonttitle=\bfseries\sffamily\small,
  coltitle=white,
  title={#1},
  arc=2mm,
  boxrule=0.8pt,
  left=9pt, right=9pt, top=7pt, bottom=7pt
}

\usepackage{booktabs}
\usepackage{longtable}
\usepackage{array}
\usepackage{calc}

\usepackage{hyperref}
\hypersetup{
  colorlinks=true,
  linkcolor=accentblue,
  urlcolor=accentblue,
  citecolor=brandblue,
  pdfborder={0 0 0}
}

\title{\Huge\bfseries\sffamily\color{brandblue} T12 --- Ereditarietà, Dynamic Binding e Polimorfismo \\[0.3em] \large\sffamily Estensione di Classi, super, Overriding vs Overloading, Polimorfismo e MavenDate v3}
\author{\textbf{Docente:} Prof. Michele Pasqua \quad---\quad \textbf{Appunti Revisione:} Wally}
\date{A.A. 2025/2026 -- Universit\`a degli Studi di Verona}

\begin{document}
\maketitle
\thispagestyle{fancy}
\vspace{-1.2em}
\noindent\textcolor{brandblue}{\rule{\textwidth}{0.8pt}}
\vspace{1.2em}

In questo blocco analizziamo in modo esaustivo i due pilastri più
importanti dell'OOP in Java: l'\textbf{Ereditarietà} e il
\textbf{Polimorfismo}
(\href{file:///home/wally/Documenti/GitHub/Programmazione-II/25-26/Teoria-e-Esercitazioni/Slides-20260902/T12\%20-\%20Inheritance\%20and\%20Polymorphism.pdf}{T12
- Inheritance and Polymorphism.pdf}). Nel codice di
\href{file:///home/wally/Documenti/GitHub/Programmazione-II/25-26/Teoria-e-Esercitazioni/Codice-20260902/Lezione12/MavenDate}{\passthrough{\lstinline!Lezione12/MavenDate!}},
la gerarchia evolve introducendo la classe astratta intermedia
\textbf{\href{file:///home/wally/Documenti/GitHub/Programmazione-II/25-26/Teoria-e-Esercitazioni/Codice-20260902/Lezione12/MavenDate/src/main/java/it/oop/core/FormattedDate.java}{\passthrough{\lstinline!FormattedDate!}}},
l'interfaccia
\textbf{\href{file:///home/wally/Documenti/GitHub/Programmazione-II/25-26/Teoria-e-Esercitazioni/Codice-20260902/Lezione12/MavenDate/src/main/java/it/oop/core/Time.java}{\passthrough{\lstinline!Time!}}},
la classe estesa
\textbf{\href{file:///home/wally/Documenti/GitHub/Programmazione-II/25-26/Teoria-e-Esercitazioni/Codice-20260902/Lezione12/MavenDate/src/main/java/it/oop/core/TimeStamp.java}{\passthrough{\lstinline!TimeStamp!}}}
e l'interfaccia standard \textbf{\passthrough{\lstinline!Comparable!}}
per l'ordinamento naturale degli array.

\begin{center}\rule{0.5\linewidth}{0.5pt}\end{center}

\subsubsection{1. Teoria: Concetti Teorici dalle Slide (Slide
T12)}\label{teoria-concetti-teorici-dalle-slide-slide-t12}

\paragraph{\texorpdfstring{A. Ereditarietà di Classe (\texttt{extends})
e Riuso (Slide
1--17)}{A. Ereditarietà di Classe (extends) e Riuso (Slide 1--17)}}\label{a.-ereditarietuxe0-di-classe-extends-e-riuso-slide-117}

\begin{itemize}
\tightlist
\item
  \textbf{Relazione ``IS-A'' vs ``HAS-A''}:

  \begin{itemize}
  \tightlist
  \item
    \emph{Ereditarietà (``IS-A'')}: un'auto \emph{è un} veicolo
    (\passthrough{\lstinline!Car extends Vehicle!}); una data italiana
    \emph{è una} data
    (\passthrough{\lstinline!ItalianDate extends Date!}).
  \item
    \emph{Composizione (``HAS-A'')}: un'auto \emph{ha un} motore
    (\passthrough{\lstinline!Car!} ha un campo
    \passthrough{\lstinline!Engine!}); una persona \emph{ha una} data di
    nascita (\passthrough{\lstinline!Person!} ha un campo
    \passthrough{\lstinline!Date!}).
  \end{itemize}
\item
  \textbf{Ereditarietà Singola in Java}:

  \begin{itemize}
  \tightlist
  \item
    In Java \textbf{non esiste l'ereditarietà multipla tra classi} (a
    differenza del C++). Una classe può estendere al massimo \textbf{una
    sola classe genitore diretta}
    (\passthrough{\lstinline!extends SuperClass!}). Questo previene il
    cosiddetto \emph{Diamond Problem}.
  \end{itemize}
\item
  \textbf{Overriding vs Overloading}:

  \begin{itemize}
  \tightlist
  \item
    \emph{Method Overloading}: metodi nella stessa classe con lo stesso
    nome ma parametri diversi (risolto staticamente a compile-time).
  \item
    \emph{Method Overriding}: una sottoclasse riscrive l'implementazione
    di un metodo ereditato mantenendo \textbf{la stessa identica
    segnatura} (nome, tipi dei parametri, ordine) e un tipo di ritorno
    identico o covariante.
  \item
    \textbf{L'Annotazione \passthrough{\lstinline!@Override!}}: Non è
    sintatticamente obbligatoria per far funzionare l'overriding, ma è
    una best practice fondamentale: dice al compilatore di verificare
    che quel metodo esista davvero nella superclasse, intercettando
    errori di battitura (es. scrivere per sbaglio
    \passthrough{\lstinline!tostring()!} con la `s' minuscola verrebbe
    segnalato come errore).
  \end{itemize}
\item
  \textbf{Costruttori ed Ereditarietà: La Parola Chiave
  \passthrough{\lstinline!super!}}:

  \begin{itemize}
  \tightlist
  \item
    Attenzione: \textbf{I costruttori NON vengono mai ereditati dalle
    sottoclassi}.
  \item
    All'interno del costruttore di una sottoclasse, la prima istruzione
    deve essere una chiamata a \passthrough{\lstinline!super(...)!}
    oppure a un altro costruttore con
    \passthrough{\lstinline!this(...)!}.
  \item
    Se il programmatore non scrive nulla, il compilatore inserisce
    automaticamente \textbf{\passthrough{\lstinline!super();!}}
    (invocazione del costruttore vuoto della superclasse).
  \item
    Attenzione: \textbf{Trappola d'esame}: Se la superclasse non
    possiede un costruttore senza parametri (perché ne ha definito uno
    con parametri e non ha aggiunto quello di default), omettere
    \passthrough{\lstinline!super(...)!} genera un \textbf{Compile-Time
    Error}!
  \item
    \passthrough{\lstinline!super.metodo()!} permette inoltre di
    invocare l'implementazione originaria della superclasse, bypassando
    l'override locale.
  \end{itemize}
\item
  \textbf{Classi e Metodi \passthrough{\lstinline!final!}}:

  \begin{itemize}
  \tightlist
  \item
    \passthrough{\lstinline!final class!}: non può essere estesa da
    nessun'altra classe (es. \passthrough{\lstinline!String!},
    \passthrough{\lstinline!Integer!}).
  \item
    \passthrough{\lstinline!final method!}: non può essere sovrascritto
    (\emph{overridden}) da nessuna sottoclasse.
  \end{itemize}
\end{itemize}

\begin{center}\rule{0.5\linewidth}{0.5pt}\end{center}

\paragraph{\texorpdfstring{B. La Classe Radice Universale:
\texttt{java.lang.Object} (Slide
18--31)}{B. La Classe Radice Universale: java.lang.Object (Slide 18--31)}}\label{b.-la-classe-radice-universale-java.lang.object-slide-1831}

\begin{itemize}
\tightlist
\item
  \emph{«The One Ring»}: in Java qualsiasi classe estende implicitamente
  \passthrough{\lstinline!java.lang.Object!}. Se una classe non
  specifica \passthrough{\lstinline!extends!}, estende direttamente
  \passthrough{\lstinline!Object!}.
\item
  \textbf{I Metodi Universali di \passthrough{\lstinline!Object!}}:

  \begin{itemize}
  \tightlist
  \item
    \passthrough{\lstinline!public String toString()!}: rappresentazione
    testuale (\passthrough{\lstinline!NomeClasse@hashcode!}).
  \item
    \passthrough{\lstinline!public boolean equals(Object obj)!}:
    confronto semantico (di default confronta l'identità con
    \passthrough{\lstinline!==!}).
  \item
    \passthrough{\lstinline!public int hashCode()!}: valore hash
    associato all'oggetto.
  \item
    \passthrough{\lstinline!public final Class<?> getClass()!}:
    restituisce il descrittore a runtime del tipo della classe.
  \end{itemize}
\item
  \textbf{Il Contratto Formale di \passthrough{\lstinline!equals!}}:
  Ogni implementazione di \passthrough{\lstinline!equals!} deve
  rispettare 5 proprietà matematiche:

  \begin{enumerate}
  \def\labelenumi{\arabic{enumi}.}
  \tightlist
  \item
    \emph{Riflessiva}: \passthrough{\lstinline!x.equals(x)!} è sempre
    \passthrough{\lstinline!true!}.
  \item
    \emph{Simmetrica}: \passthrough{\lstinline!x.equals(y)!} ritorna
    \passthrough{\lstinline!true!} se e solo se
    \passthrough{\lstinline!y.equals(x)!} ritorna
    \passthrough{\lstinline!true!}.
  \item
    \emph{Transitiva}: se \passthrough{\lstinline!x.equals(y)!} è
    \passthrough{\lstinline!true!} e
    \passthrough{\lstinline!y.equals(z)!} è
    \passthrough{\lstinline!true!}, allora
    \passthrough{\lstinline!x.equals(z)!} è
    \passthrough{\lstinline!true!}.
  \item
    \emph{Consistente}: invocazioni multiple producono lo stesso
    risultato finché lo stato non muta.
  \item
    \emph{Non-Nullità}: per qualsiasi
    \passthrough{\lstinline"x != null"}, la chiamata
    \passthrough{\lstinline!x.equals(null)!} \textbf{deve restituire
    \passthrough{\lstinline!false!}} senza lanciare eccezioni.
  \end{enumerate}
\end{itemize}

\begin{center}\rule{0.5\linewidth}{0.5pt}\end{center}

\paragraph{C. Polimorfismo e Dynamic Method Dispatch (Late Binding)
(Slide
32--50)}\label{c.-polimorfismo-e-dynamic-method-dispatch-late-binding-slide-3250}

\begin{itemize}
\item
  \textbf{Principio di Sostituzione di Liskov (LSP)}: Se \(S\) è
  sottotipo di \(T\), qualsiasi porzione di codice scritta per
  interagire con \(T\) deve poter operare in modo trasparente e corretto
  con un'istanza di \(S\).
\item
  \textbf{Tipo Statico vs Tipo Dinamico}:

\begin{lstlisting}[language=Java]
Date d = new ItalianDate(1, 1, 1970);
\end{lstlisting}

  \begin{itemize}
  \tightlist
  \item
    \textbf{Tipo Statico (a Compile-Time)}: è il tipo dichiarato della
    variabile reference (\passthrough{\lstinline!Date!}). Il compilatore
    accetta solo invocazioni a metodi e accessi a campi definiti nella
    classe \passthrough{\lstinline!Date!}.
  \item
    \textbf{Tipo Dinamico (a Runtime)}: è il tipo effettivo dell'oggetto
    allocato nello Heap (\passthrough{\lstinline!ItalianDate!}).
  \end{itemize}
\item
  \textbf{Dynamic Method Dispatch (Late Binding)}: Quando a runtime
  viene eseguita una chiamata a un metodo d'istanza (es.
  \passthrough{\lstinline!d.toString()!}), la JVM consulta la
  \textbf{Virtual Method Table (vtable)} dell'oggetto concreto nello
  Heap e seleziona \textbf{la versione del metodo definita dal Tipo
  Dinamico}! Non importa che \passthrough{\lstinline!d!} sia dichiarata
  come \passthrough{\lstinline!Date!}: verrà eseguito il
  \passthrough{\lstinline!toString()!} specializzato di
  \passthrough{\lstinline!ItalianDate!}.
\item
  \textbf{Regole di Conversione di Tipo (Casting tra Oggetti)}:

  \begin{itemize}
  \item
    \textbf{Upcasting (verso la superclasse)}:
    \passthrough{\lstinline!Date d = new ItalianDate(...);!}
    \(\implies\) \textbf{Implicito, automatico e sempre sicuro al
    100\%}.
  \item
    \textbf{Downcasting (verso la sottoclasse)}:
    \passthrough{\lstinline!ItalianDate id = (ItalianDate) d;!}
    \(\implies\) \textbf{Esplicito obbligatorio}.

    \begin{itemize}
    \tightlist
    \item
      Se l'oggetto concreto nello Heap è effettivamente un'istanza
      compatibile, il cast ha successo.
    \item
      Se l'oggetto concreto nello Heap è di tipo diverso (es. un
      \passthrough{\lstinline!Date!} generico o un
      \passthrough{\lstinline!AmericanDate!}), la JVM interrompe
      l'esecuzione e lancia una
      \textbf{\passthrough{\lstinline!ClassCastException!}}!
    \end{itemize}
  \item
    \textbf{Downcasting Sicuro con
    \passthrough{\lstinline!instanceof!}}: Per evitare crash a runtime
    si verifica preventivamente la compatibilità:

\begin{lstlisting}[language=Java]
if (d instanceof ItalianDate) {
    ItalianDate id = (ItalianDate) d;
}
\end{lstlisting}
  \end{itemize}
\end{itemize}

\begin{center}\rule{0.5\linewidth}{0.5pt}\end{center}

\subsubsection{\texorpdfstring{2. Codice: Evoluzione del Codice:
\href{file:///home/wally/Documenti/GitHub/Programmazione-II/25-26/Teoria-e-Esercitazioni/Codice-20260902/Lezione12/MavenDate}{\texttt{Lezione12/MavenDate}}}{2. Codice: Evoluzione del Codice: Lezione12/MavenDate}}\label{codice-evoluzione-del-codice-lezione12mavendate}

In \passthrough{\lstinline!Lezione12!} la gerarchia del progetto diventa
ricca e completa:

\begin{lstlisting}
classDiagram
    class Object {
        +toString()
        +equals(Object)
    }
    class Comparable {
        <<interface>>
        +compareTo(Object)
    }
    class Time {
        <<interface>>
        +getSeconds()
        +getMinutes()
        +getHours()
    }
    class Date {
        #day: int
        -month: int
        -year: int
        +toString()
        +equals(Object)
        +compareTo(Object)
    }
    class FormattedDate {
        <<abstract>>
        #format: String
        #months: String[]
        +printFormat() final
        +getMonthAsString() final
        +prettyPrint()* String
    }
    class ItalianDate {
        +prettyPrint()
        +toString()
    }
    class AmericanDate {
        +prettyPrint()
        +toString()
    }
    class TimeStamp {
        #seconds: int
        #minutes: int
        #hours: int
        +toString()
        +equals(Object)
    }

    Object <|-- Date
    Comparable <|.. Date
    Date <|-- FormattedDate
    FormattedDate <|-- ItalianDate
    FormattedDate <|-- AmericanDate
    Date <|-- TimeStamp
    Time <|.. TimeStamp
\end{lstlisting}

\paragraph{\texorpdfstring{1. L'Interfaccia \texttt{Comparable} in
\href{file:///home/wally/Documenti/GitHub/Programmazione-II/25-26/Teoria-e-Esercitazioni/Codice-20260902/Lezione12/MavenDate/src/main/java/it/oop/core/Date.java}{\texttt{Date.java}}}{1. L'Interfaccia Comparable in Date.java}}\label{linterfaccia-comparable-in-date.java}

\passthrough{\lstinline!Date!} implementa l'interfaccia standard
\passthrough{\lstinline!java.lang.Comparable!} per consentire
l'ordinamento naturale:

\begin{lstlisting}[language=Java]
public class Date implements Comparable {
    ...
    @Override
    public int compareTo(Object other) {
        Date otherAsDate = (Date) other;
        int diff = this.year - otherAsDate.getYear();
        if (diff != 0) return diff;
        diff = this.month - otherAsDate.getMonth();
        if (diff != 0) return diff;
        return this.day - otherAsDate.getDay();
    }
}
\end{lstlisting}

\emph{Logica di confronto}: confronta prima gli anni, se uguali
confronta i mesi, se uguali confronta i giorni. Restituisce un intero
negativo se \passthrough{\lstinline!this < other!}, zero se uguali,
positivo se \passthrough{\lstinline!this > other!}.

\begin{center}\rule{0.5\linewidth}{0.5pt}\end{center}

\paragraph{\texorpdfstring{2. L'Astrazione con Classe Astratta:
\href{file:///home/wally/Documenti/GitHub/Programmazione-II/25-26/Teoria-e-Esercitazioni/Codice-20260902/Lezione12/MavenDate/src/main/java/it/oop/core/FormattedDate.java}{\texttt{FormattedDate.java}}}{2. L'Astrazione con Classe Astratta: FormattedDate.java}}\label{lastrazione-con-classe-astratta-formatteddate.java}

Centralizza i dati di formattazione evitando duplicazioni tra
\passthrough{\lstinline!ItalianDate!} e
\passthrough{\lstinline!AmericanDate!}:

\begin{lstlisting}[language=Java]
public abstract class FormattedDate extends Date {
    protected final String format;
    protected final String[] months;

    public FormattedDate(int day, int month, int year, String format, String[] months) {
        super(day, month, year);
        this.format = format;
        this.months = months;
    }

    // Metodi final: le sottoclassi non possono sovrascriverli (Template Method Pattern)
    public final String printFormat() { return format; }
    public final String getMonthAsString() { return months[getMonth() - 1]; }

    // Metodo astratto: DEVE essere implementato dalle sottoclassi concrete
    public abstract String prettyPrint();
}
\end{lstlisting}

\begin{center}\rule{0.5\linewidth}{0.5pt}\end{center}

\paragraph{\texorpdfstring{3. Ereditarietà Multipla di Tipo:
\href{file:///home/wally/Documenti/GitHub/Programmazione-II/25-26/Teoria-e-Esercitazioni/Codice-20260902/Lezione12/MavenDate/src/main/java/it/oop/core/TimeStamp.java}{\texttt{TimeStamp.java}}}{3. Ereditarietà Multipla di Tipo: TimeStamp.java}}\label{ereditarietuxe0-multipla-di-tipo-timestamp.java}

Mostra come combinare l'estensione di una classe concreta con
l'implementazione di un'interfaccia:

\begin{lstlisting}[language=Java]
public class TimeStamp extends Date implements Time {
    protected final int seconds;
    protected final int minutes;
    protected final int hours;

    public TimeStamp(int seconds, int minutes, int hours, int day, int month, int year) {
        super(day, month, year);
        this.seconds = seconds;
        this.minutes = minutes;
        this.hours = hours;
    }

    @Override
    public String toString() {
        // Riusa il toString() della superclasse aggiungendo l'orario con zero-padding a due cifre:
        return String.format("%s[%02d:%02d:%02d]", super.toString(), hours, minutes, seconds);
    }

    @Override
    public boolean equals(Object other) {
        // 1. Verifica prima l'uguaglianza della data con super.equals():
        if (super.equals(other) == false) return false;
        if (!(other instanceof TimeStamp)) return false;
        TimeStamp otherAsTimeStamp = (TimeStamp) other;
        // 2. Poi verifica l'orario:
        return hours == otherAsTimeStamp.getHours() &&
               minutes == otherAsTimeStamp.getMinutes() &&
               seconds == otherAsTimeStamp.getSeconds();
    }
}
\end{lstlisting}

\begin{center}\rule{0.5\linewidth}{0.5pt}\end{center}

\paragraph{\texorpdfstring{4. Dimostrazione Pratica in
\href{file:///home/wally/Documenti/GitHub/Programmazione-II/25-26/Teoria-e-Esercitazioni/Codice-20260902/Lezione12/MavenDate/src/main/java/it/oop/ui/MainDate.java}{\texttt{MainDate.java}}}{4. Dimostrazione Pratica in MainDate.java}}\label{dimostrazione-pratica-in-maindate.java}

\begin{lstlisting}[language=Java]
Time ts1 = new TimeStamp(0, 0, 10, 10, 11, 2025);
System.out.println(ts1.toString()); // Esegue TimeStamp.toString() tramite Late Binding!

// ts1.getDay(); // COMPILE-TIME ERROR: il tipo statico Time non possiede il metodo getDay()!
System.out.println(((Date) ts1).getDay()); // OK: cast esplicito a Date!

Date[] arr = { ts2, itDate, new Date(9, 9, 2025) };
Arrays.sort(arr); // Ordina l'array eterogeneo grazie a compareTo()!
System.out.println(Arrays.toString(arr));
\end{lstlisting}

\textbf{Output a video prodotto dall'ordinamento polimorfico}:

\begin{lstlisting}
[y2025m9d9, 3/11/2025, y2025m11d10[10:00:00]]
\end{lstlisting}

\begin{itemize}
\tightlist
\item
  L'ordinamento colloca correttamente prima il 9 settembre, poi il 3
  novembre (\passthrough{\lstinline!itDate!}) e infine il 10 novembre
  (\passthrough{\lstinline!TimeStamp!}).
\item
  In fase di stampa, ciascun elemento viene formattato dal
  \textbf{proprio} metodo \passthrough{\lstinline!toString()!} dinamico.
\end{itemize}

\begin{center}\rule{0.5\linewidth}{0.5pt}\end{center}

\begin{studynote}[Nota di Studio] Con la \textbf{Lezione 12} abbiamo analizzato la gerarchia
completa di ereditarietà, l'overriding con
\passthrough{\lstinline!super!}, il late binding, la classe astratta
\passthrough{\lstinline!FormattedDate!} e l'ordinamento con
\passthrough{\lstinline!Comparable!}.

Il prossimo blocco è la \textbf{Lezione 13 / Slide T13: \emph{Abstract
Classes and Interfaces}}: - Differenze ontologiche e strutturali tra
\textbf{Classi Astratte} e \textbf{Interfacce}. - Metodi
\passthrough{\lstinline!default!} e metodi
\passthrough{\lstinline!static!} nelle interfacce (Java 8+). -
Ereditarietà multipla tramite interfacce e risoluzione dei conflitti
(\emph{Diamond Problem} sui metodi default). - Evoluzione di
\passthrough{\lstinline!MavenDate!} in
\passthrough{\lstinline!Lezione13!}: - La classe
\passthrough{\lstinline!Date!} viene ripensata per gestire il calendario
completo. - Analisi delle nuove interfacce e classi di supporto.
\end{studynote}

\textbf{Dammi conferma per aprire la Lezione 13 / T13 e analizzare a
fondo Classi Astratte e Interfacce!}


\end{document}
